\subsection{Module 3: Ras Activation, Adaptation and Amplification}\label{sec:m3}

\paragraph{Role.}
The earliest polarised response downstream of the G protein is the activation of Ras.
RasG and RasC are \mk{M3-trans}{activated rapidly and transiently by cAMP in aggregation-competent cells}~\cite{kae2004}, and \mk{M3-pi3k}{RasG is the upstream activator of PI3K}~\cite{sasaki2004,takeda2012}.
Adaptation takes place between the G protein and Ras: \mk{M3-g-persist}{G-protein activation does not decline during continuous stimulation, whereas the downstream responses subside}~\cite{janetopoulos2001,elzie2009}, and \mk{M3-g-steps}{under two successive cAMP steps the G-protein dissociation rises step-like and persists, while Ras responds transiently to each step}~\cite{xu2022}.
\mk{M3-adapt}{The activation of RasG shows near-perfect adaptation over a wide range of concentrations}~\cite{takeda2012}; among the two simple topologies that can produce perfect adaptation, \mk{M3-iffl}{only an incoherent feed-forward loop fitted the measured responses}~\cite{takeda2012}, in which \mk{M3-prop}{the activator drives an excitatory and an inhibitory node proportionally}~\cite{goentoro2009,takeda2012}, \mk{M3-gap-glob}{possibly with a RasGAP as the global inhibitor}~\cite{takeda2012}.
An incoherent feed-forward loop of this type has the structure of a local-excitation--global-inhibition (LEGI) gradient sensor~\cite{levchenko2002,ma2004}.
In the Ras network, \mk{M3-legi}{the activator and the inactivator of Ras, a RasGEF and a RasGAP, are both activated by receptor and G-protein signalling}~\cite{takeda2012,meierschellersheim2006}, and such a loop can provide \mk{M3-fcd3}{fold-change detection}~\cite{goentoro2009}.
The RasGAP NF1 (DdNF1) is \mk{M3-nf1}{a major regulator of Ras activity and its loss makes Ras activity spatially and temporally unregulated}~\cite{zhang2008}; \mk{M3-nf1-rasg}{in its absence the activation of RasG is delayed and extended, whereas RasD, Rap1 and RasC are unaffected}~\cite{zhang2008}.
\mk{M3-rect}{Ras activation is suppressed when the cAMP concentration falls (``rectification''), which arises naturally in an incoherent feed-forward loop with zero-order ultrasensitivity}~\cite{nakajima2014}.
\mk{M3-ras-no-down}{The Ras response to uniform cAMP needs none of the downstream PI3K, TorC2, PLA2 and sGC pathways and is unchanged when all four are inhibited, also together with latrunculin~A and in a pkbR1-null background}~\cite{kortholt2011}.
Finally, the Ras/PI3K/PIP3 network is \mk{M3-excit}{excitable: it fires without external cues}~\cite{devreotes2017,fukushima2019,miao2017}, \mk{M3-bias}{a stimulus biases threshold and timing}~\cite{fukushima2019,miao2017,devreotes2017,iglesias2012}, and \mk{M3-excit-intr}{Ras remains excitable without the PIP3, TorC2, PLA2 and sGC pathways and without a functional actin cytoskeleton}~\cite{fukushima2019,iwamoto2025}.
\mk{M3-dnf}{Excitability requires a positive feedback for the all-or-none response and a delayed negative feedback that terminates the response and is followed by a refractory period}~\cite{fukushima2019}.
Module~3 therefore combines (i)~a fast G$\beta\gamma$-driven excitor and a slower G$\beta\gamma$-driven inhibitor (adaptation and spatial comparison), (ii)~a saturable GAP (rectification), (iii)~a positive feedback from PIP3 onto RasG (amplification) and (iv)~a delayed negative feedback read from RasG-GTP, the brake (termination of excursions that the stimulus did not cause; proposed).
An integral-feedback (antithetic) arm that was tested earlier was removed, because \mk{M3-noint}{integral feedback gives oscillations and stimulus-dependent kinetics that are inconsistent with the measured Ras response}~\cite{takeda2012}.

\paragraph{Adaptation and brake.}
The two negative arms act on the same RasG-GTP but read different variables, and this is what separates them.
The inhibitor of the incoherent feed-forward loop (3.4e, 3.5, 3.7z) reads the input, G$\beta\gamma$, as the excitor does.
It sets the adapted level, which is therefore independent of the dose, and it alone acts when the stimulus falls: \mk{M3-undershoot}{after a sudden decrease of cAMP the Ras reporter returns rapidly to the cytosol and then slowly to its basal amount}~\cite{takeda2012}, i.e.\ RasG-GTP transiently falls below its resting level because the slow inhibitor outlasts the fast excitor.
The brake (3.3c$_\mathrm{R}$, 3.3d, 3.3e) reads the output, RasG-GTP.
Its purpose is the delayed negative feedback that excitability requires (see Role): it is the only arm that can terminate a Ras excursion that the stimulus did not cause, such as a spontaneous firing of the excitable network, because the inhibitor follows G$\beta\gamma$ and is not raised by such an excursion.
After a brief stimulus both arms contribute to the refractory period, since B$^*$ and the inhibitor both decay with $\tau=10$\,s; the measured target is that \mk{M3-refr}{with paired 2-s mechanical stimuli and an F-actin reporter, the second stimulus elicits no response at intervals below 12\,s, responsiveness then recovers with a half-time of about 7\,s as observed for chemoattractant, and a refractory period is one of the features of excitable systems}~\cite{artemenko2016}.
The brake is constructed so that it does not set the adapted level (design choice): B$^*$ decays in first order and has no set point, and its half-activation lies at the driven RasG-GTP level, so that its removal flux, which is proportional to $[\mathrm{B}^*][\mathrm{RasG^{GTP}}]$ and hence roughly to the square of RasG-GTP while B is far from saturation, is small at rest and in the adapted plateau and large only at the peak.
It is not the integral controller that was rejected for Ras: in that model \mk{M3-int-def}{the GAP is produced in proportion to Ras-GTP and removed at a constant, zero-order rate}~\cite[Suppl.\ Information]{takeda2012}, which holds Ras-GTP at a set point, whereas B$^*$ is removed in first order.
That the brake does not adapt is a deviation from the literature, in which a feedback read from Ras is proposed as an adaptation mechanism: \mk{M3-nfblb-def}{a negative feedback loop with a buffer node, in which the output is shut down by an inhibitor induced by the output itself, is one of the two simple topologies that produce adaptation}~\cite{xu2022}, and \mk{M3-nfblb-fit}{models in which the RasGAP is activated by G$\alpha_2$-GTP, by Ras-GTP or by both reproduce transient, adaptive Ras signalling}~\cite{xu2022}.
The model assigns adaptation to the incoherent feed-forward loop, the only topology that fitted the measured responses in the comparison with integral feedback (see Role), and keeps the Ras-read feedback out of the adapted state; this division of labour is a design choice and is not shown by the sources.
The absence of a second excursion after a step is a test of the tuning.
During a sustained cAMP step both arms act; they are distinguished by the undershoot after removal of the stimulus (inhibitor only) and by the termination of spontaneous excursions (brake only).

\begin{table*}[!t]
\centering
\caption{Module~3 reactions in the default configuration. GEF$_R$ and GAP are the cytosolic RasG exchange factor and GTPase-activating protein (asterisk: active form); $\mathrm{RasG}^{\mathrm{GTP/GDP}}$ are the membrane-bound Ras states; B is the substrate of the delayed brake (3.3c$_\mathrm{R}$--e), activated by RasG-GTP. 3.1--3.7z form the adaptation (incoherent feed-forward loop, read from G$\beta\gamma$), 3.3b the amplification and 3.3c$_\mathrm{R}$--e the brake (read from RasG-GTP). Reactions marked $^\dagger$ are \emph{proposed}: the mechanism is a candidate, and they are kept only if tuning shows that they are needed (see text). A coloured reaction is the one to which the equally coloured statement in the text and the equally coloured passage in the cited paper refer; a small square appends a further statement.}
\label{tab:m3rxn}
\footnotesize
\renewcommand{\arraystretch}{1.15}
\begin{tabularx}{\textwidth}{@{}l>{\raggedright\arraybackslash}p{0.37\textwidth}>{\raggedright\arraybackslash}X@{}}
\toprule
ID & Reaction & Propensity (rate law)\\
\midrule
\mkn{M3-kort}{M3-prop,M3-legi}{3.1} & \mkn{M3-kort}{M3-prop,M3-legi}{$\mathrm{G}_{\beta\gamma}+\mathrm{GEF_R} \to \mathrm{GEF_R^*}+\mathrm{G}_{\beta\gamma}$} & $k_{E,\mathrm{on}}[\mathrm{G}_{\beta\gamma}][\mathrm{GEF_R}]$ \\
\mkn{M3-peak-fast}{M3-tau-fit}{3.2} & \mkn{M3-peak-fast}{M3-tau-fit}{$\mathrm{GEF_R^*} \to \mathrm{GEF_R}$} & $k_{E,\mathrm{off}}[\mathrm{GEF_R^*}]$ \\
\mkn{M3-gefr}{M3-trans,M3-gef-exch}{3.3} & \mkn{M3-gefr}{M3-trans,M3-gef-exch}{$\mathrm{GEF_R^*}+\mathrm{RasG}^{\mathrm{GDP}} \to \mathrm{RasG}^{\mathrm{GTP}}+\mathrm{GEF_R^*}$} & $k_{R,\mathrm{on}}[\mathrm{GEF_R^*}][\mathrm{RasG}^{\mathrm{GDP}}]$ \\
\mkn{M3-prop}{M3-gap-glob,M3-legi,M3-iffl,M3-gap-gbg}{3.4e} & \mkn{M3-prop}{M3-gap-glob,M3-legi,M3-iffl,M3-gap-gbg}{$\mathrm{G}_{\beta\gamma}+\mathrm{GAP} \to \mathrm{GAP^*}+\mathrm{G}_{\beta\gamma}$} & $k_{I,\mathrm{on}}[\mathrm{G}_{\beta\gamma}][\mathrm{GAP}]$ \\
\mkt{M3-basal-gap}{3.4b} & \mkt{M3-basal-gap}{$\mathrm{GAP} \to \mathrm{GAP^*}$} & $k_{I,\mathrm{b}}[\mathrm{GAP}]$ \\
\mkn{M3-tau-fit}{M3-slower-inh,M3-undershoot}{3.5} & \mkn{M3-tau-fit}{M3-slower-inh,M3-undershoot}{$\mathrm{GAP^*} \to \mathrm{GAP}$} & $k_{I,\mathrm{off}}[\mathrm{GAP^*}]$ \\
\mkn{M3-rect}{M3-nf1kin,M3-nf1-rasg}{3.7z} & \mkn{M3-rect}{M3-nf1kin,M3-nf1-rasg}{$\mathrm{GAP^*}+\mathrm{RasG}^{\mathrm{GTP}} \to \mathrm{RasG}^{\mathrm{GDP}}+\mathrm{GAP^*}$} & $k_\mathrm{cat}[\mathrm{GAP^*}]\,\dfrac{[\mathrm{RasG}^{\mathrm{GTP}}]}{K_\mathrm{NF1}+[\mathrm{RasG}^{\mathrm{GTP}}]}$ \\
\mkn{M3-pfb}{M3-pip3-gef,M3-ly-red,M3-pip3-small}{3.3b} & \mkn{M3-pfb}{M3-pip3-gef,M3-ly-red,M3-pip3-small}{$\mathrm{RasG}^{\mathrm{GDP}}+\mathrm{PIP3} \to \mathrm{RasG}^{\mathrm{GTP}}+\mathrm{PIP3}$} & $k_\mathrm{PIP3,R}[\mathrm{PIP3}][\mathrm{RasG}^{\mathrm{GDP}}]$ \\
\mkn{M3-nfblb-def}{M3-c2gap-ras,M3-dnf}{3.3c$_\mathrm{R}^\dagger$} & \mkn{M3-nfblb-def}{M3-c2gap-ras,M3-dnf}{$\mathrm{RasG}^{\mathrm{GTP}}+\mathrm{B} \to \mathrm{B}^*+\mathrm{RasG}^{\mathrm{GTP}}$} & $k_{B,\mathrm{on}}[\mathrm{RasG}^{\mathrm{GTP}}][\mathrm{B}]$ \\
\mkn{M3-refr}{M3-slow-buf}{3.3d$^\dagger$} & \mkn{M3-refr}{M3-slow-buf}{$\mathrm{B}^* \to \mathrm{B}$} & $k_{B,\mathrm{off}}[\mathrm{B}^*]$ \\
\mkn{M3-c2gap-ras}{M3-nfblb-fit}{3.3e$^\dagger$} & \mkn{M3-c2gap-ras}{M3-nfblb-fit}{$\mathrm{B}^*+\mathrm{RasG}^{\mathrm{GTP}} \to \mathrm{RasG}^{\mathrm{GDP}}+\mathrm{B}^*$} & $k_{B,\mathrm{hyd}}[\mathrm{B}^*][\mathrm{RasG}^{\mathrm{GTP}}]$ \\
\bottomrule
\end{tabularx}
\end{table*}

\begin{table*}[!t]
\centering
\caption{Module~3 constants. Rate constants and pool sizes are paired: $k_{E,\mathrm{on}}[\mathrm{GEF_R}]_\mathrm{tot}$, $k_{R,\mathrm{on}}[\mathrm{GEF_R}]_\mathrm{tot}$ and $k_{I,\mathrm{on}}[\mathrm{GAP}]_\mathrm{tot}$ were held fixed when the pools were enlarged, so that only shot noise changed.}
\label{tab:m3const}
\footnotesize
\renewcommand{\arraystretch}{1.15}
\begin{tabularx}{\textwidth}{@{}l r l c >{\raggedright\arraybackslash}X@{}}
\toprule
Symbol & Value & Unit & Basis & Meaning, justification, source\\
\midrule
$[\mathrm{RasG}]_\mathrm{tot}$ & \num{38500} & $\mathrm{molec/voxel}$ & E & $5\times10^5$ per cell; large for low shot noise. The copy number is not well measured. \\
$[\mathrm{GEF_R}]_\mathrm{tot}$ & \num{138600} & $\mathrm{molec/voxel}$ & E & Excitor pool ($1.8\times10^6$ per cell). \mkt{M3-gefr}{RasGEFR activates RasG}~\cite{kae2007}. \\
$[\mathrm{GAP}]_\mathrm{tot}$ & \num{46200} & $\mathrm{molec/voxel}$ & E & Inhibitor pool ($6\times10^5$ per cell); \mkt{M3-nf1}{a RasGAP (NF1) is the global inhibitor}~\cite{zhang2008}. \\
$k_{E,\mathrm{on}}$ & \num{3.8e-3} & $\mu\mathrm{M}^{-1}\mathrm{s}^{-1}$ & C & G$\beta\gamma$-driven GEF activation; \mkt{M3-kort}{the initial Ras response requires G$\beta$}~\cite{kortholt2013}. \\
$k_{E,\mathrm{off}}$ & \num{0.5} & $\mathrm{s^{-1}}$ & C & GEF inactivation, $\tau_E=2$\,s: \mkt{M3-peak-fast}{the excitor is fast, so the Ras peak occurs within seconds}~\cite{kae2004,takeda2012}; close to \mkt{M3-tau-fit}{the fitted $0.4\,\mathrm{s^{-1}}$}~\cite{takeda2012}. \\
$k_{R,\mathrm{on}}$ & \num{0.333} & $\mu\mathrm{M}^{-1}\mathrm{s}^{-1}$ & C & GEF-catalysed GDP$\to$GTP exchange on RasG. \\
$k_{I,\mathrm{on}}$ & \num{7.67e-4} & $\mu\mathrm{M}^{-1}\mathrm{s}^{-1}$ & D & G$\beta\gamma$-driven GAP activation. It reads the \emph{same} input as 3.1, as \mkt{M3-prop}{an incoherent feed-forward loop requires}~\cite{takeda2012,goentoro2009}. Calibrated value $2.53\times10^{-4}$ multiplied by $3.03$ together with $k_{I,\mathrm{off}}$, which keeps the activated fraction; \mkt{M3-tau-pair}{the fit tied activation and deactivation rates in the same way}~\cite{takeda2012}. \\
$k_{I,\mathrm{off}}$ & \num{0.1} & $\mathrm{s^{-1}}$ & L & GAP inactivation, $\tau_I=10$\,s, \mkt{M3-tau-fit}{the fitted value, which sets the time of the return to the basal amount}~\cite{takeda2012} (until 2026-10: $0.033\,\mathrm{s^{-1}}$); \mkt{M3-slower-inh}{the inhibitor is slower than the excitor, which produces adaptation}~\cite{takeda2012}. Not yet recalibrated against the downstream modules. \\
$k_{I,\mathrm{b}}$ & \num{1.52e-3} & $\mathrm{s^{-1}}$ & D & Stimulus-independent GAP activity. Without it the resting Ras/PIP3 loop latches; consistent with \mkt{M3-basal-gap}{a basal RasGAP}~\cite{zhang2008,xu2021c2gap}. Calibrated as the minimum value giving the low rest state ($5\times10^{-4}$ at $\tau_I=30$\,s), multiplied by $3.03$ with $k_{I,\mathrm{off}}$, which keeps the basal activated fraction. \\
$K_\mathrm{NF1}$ & \num{0.2} & $\mu\mathrm{M}$ & L & \mkt{M3-nf1kin}{Michaelis constant of the NF1 GAP domain, $0.3\,\mu$M}~\cite{wiesmuller1992,coyle2016}; the value is rounded to $0.2$. Mammalian catalytic domain; no \emph{Dictyostelium} kinetics exist. \\
$k_\mathrm{cat}$ & \num{15.8} & $\mathrm{s^{-1}}$ & D & Derived from the resting distance to the zero-order switch, $\theta_\mathrm{rest}=0.5$. Only the product $k_\mathrm{cat}[\mathrm{GAP}^*]$ enters; it is 11$\times$ the mammalian \mkt{M3-nf1kin}{$k_\mathrm{cat}=1.4$\,s$^{-1}$}~\cite{coyle2016,wiesmuller1992}, so the active-GAP pool stands for more enzyme than is measured. \\
$k_\mathrm{PIP3,R}$ & \num{0.04} & $\mu\mathrm{M}^{-1}\mathrm{s}^{-1}$ & C & PIP3$\to$RasG feedback (3.3b), \mkt{M3-pip3-gef}{the PIP3 term of the Ras GEF reaction}~\cite{fukushima2019}. The value committed for this form in an earlier calibration (higher gains raised the variance but not the contrast, see text); not re-tuned after the changes of 2026-10. Kept small, since \mkt{M3-ly-red}{PI3K inhibition reduces but does not abolish Ras activation}~\cite{sasaki2004}. Expressed relative to the legacy PIP2 pool (Sec.~\ref{sec:m4}). \\
$k_{B,\mathrm{on}}$ & \num{0.0753} & $\mu\mathrm{M}^{-1}\mathrm{s}^{-1}$ & E & Brake activation by RasG-GTP (3.3c$_\mathrm{R}$); half-activation at 6400 RasG-GTP molecules/voxel, the driven level recorded in the code (design anchor, not tuned). Proposed brake. \\
$k_{B,\mathrm{off}}$ & \num{0.1} & $\mathrm{s^{-1}}$ & D & Brake decay ($\tau=10$\,s), from \mkt{M3-refr}{the recovery half-time of about 7\,s after a refractory period}~\cite{artemenko2016} as $\tau=t_{1/2}/\ln2$, assuming that B$^*$ governs the recovery; measured with an F-actin reporter after mechanical stimulation, not for Ras. \\
$k_{B,\mathrm{hyd}}$ & \num{39.8} & $\mu\mathrm{M}^{-1}\mathrm{s}^{-1}$ & D & Chosen so that at half activation the brake removes RasG-GTP at $0.6\times$ the rate of the LEGI GAP arm at the driven state (design target, not tuned). Proposed brake. \\
$[\mathrm{B}]_\mathrm{tot}$ & \num{1540} & $\mathrm{molec/voxel}$ & E & Brake pool; the mean braking action is independent of pool size, so it was raised to reduce noise. \\
$D_{\mathrm{GEF_R}},D_{\mathrm{RasG}}$ & \num{0.1} & $\mu\mathrm{m^2/s}$ & E & Local excitor and membrane-bound Ras (prenylated anchor). \\
$D_{\mathrm{GAP}},D_\mathrm{B}$ & \num{20} & $\mu\mathrm{m^2/s}$ & E & Global inhibitor and global brake. Cytosolic value; for B$^*$, a membrane-attached species in the code, it contradicts the local purpose of the brake (open, see 3.3c$_\mathrm{R}$--e); $0.4\,\mu\mathrm{m^2/s}$ gives the local brake. \\
\bottomrule
\end{tabularx}
\end{table*}


\paragraph{Reactions.}
\begin{description}\itemsep0pt
\item[3.1/3.2] The excitor GEF$_\mathrm{R}$ is activated by free G$\beta\gamma$ and inactivated with $\tau_E=2$\,s. \mk{M3-gefr}{RasGEFR is the exchange factor of RasG}~\cite{kae2007}, and \mk{M3-kort}{the initial Ras response depends on G$\beta$}~\cite{kortholt2013}. The short relaxation time makes the RasG peak follow the input within seconds. It is close to the fitted value: \mk{M3-tau-fit}{in the fitted incoherent feed-forward model the GEF is inactivated at $0.4\,\mathrm{s^{-1}}$ and the GAP at $0.1\,\mathrm{s^{-1}}$, the GAP rate sets the time of the return to the basal amount, and the GEF rate must exceed the GAP rate}~\cite[Suppl.\ Information]{takeda2012}.
\item[3.3] The active GEF exchanges GDP for GTP on membrane-bound RasG, which produces RasG-GTP. This is the general mechanism of Ras activation: \mk{M3-gef-exch}{Ras GEFs catalyse the exchange of GDP for GTP and thereby activate Ras}~\cite{kae2004,kae2007,xu2022}.
\item[3.4e/3.4b/3.5] The inhibitor GAP is activated by G$\beta\gamma$ (3.4e), i.e.\ by the same variable as the excitor, and decays with $\tau_I=10$\,s (3.5), the GAP rate of the fit quoted under 3.1/3.2 (until 2026-10 the model used $\tau_I=30$\,s). A G$\beta\gamma$-armed GAP is also the choice of an earlier model, in which \mk{M3-gap-gbg}{a RasGAP is activated by G$\beta\gamma$, translocates to the membrane and deactivates Ras}~\cite{meierschellersheim2006}. The requirement that both arms read the same input is essential: if they read different upstream variables with different dose dependence, the ratio of excitation to inhibition acquires a dose dependence of its own and the stage inverts instead of adapting (this was found when the inhibitor read the receptor occupancy). A small basal activation (3.4b) represents the stimulus-independent GAP at rest, which is needed because the resting Ras/PIP3 loop otherwise latches on; a basal RasGAP is consistent with \mk{M3-basal-gap}{the enhanced basal Ras activity of GAP-deficient cells}~\cite{zhang2008,xu2021c2gap}.
\item[3.7z] The active GAP hydrolyses RasG-GTP with Michaelis--Menten kinetics. Because $K_\mathrm{NF1}$ is far below the Ras pool, the enzyme operates in the zero-order regime, where a graded input is turned into a switch-like output and \mk{M3-rect}{the response to a falling stimulus is suppressed}~\cite{nakajima2014}. \mk{M3-nf1kin}{The Michaelis constant of the NF1 GAP domain is $0.3\,\mu$M and its $k_\mathrm{cat}$ is $1.4$\,s$^{-1}$ (mammalian catalytic domain)}~\cite{wiesmuller1992,coyle2016}. The plain mass-action form of this step is available and is not used.
\item[3.3b] PIP3 converts membrane-bound RasG-GDP into RasG-GTP. Together with 4.1 (RasG-GTP recruits PI3K) and 4.3 (PI3K makes PIP3) this closes the \mk{M3-pfb}{Ras$\to$PI3K$\to$PIP3$\to$Ras positive-feedback loop, which is the amplifier of the excitable network}~\cite{fukushima2019}, consistent with \mk{M3-kat}{the amplification between the G protein and Ras}~\cite{kataria2013}. The form is that of \mk{M3-pip3-gef}{the published Ras/PIP3 model, whose Ras GEF reaction has a basal term and a term for the feedback from PIP3}~\cite{fukushima2019}. The feedback contributes but is not required: \mk{M3-ly-red}{initial Ras activation requires neither PI3K activity nor F-actin, but its level is reduced by LY294002 or latrunculin}~\cite{sasaki2004}, \mk{M3-ras-pip3-indep}{PIP3 production and degradation are not necessary for Ras wave generation}~\cite{fukushima2019}, and \mk{M3-pip3-small}{elevated Ras activity depends largely on decreased PI(3,4)P$_2$, with a contribution of feedback from PIP3}~\cite{li2018}.
The molecule on which PIP3 acts is not identified; a route through GAPs is possible, since \mk{M3-li-gap}{RasGAP2 and RapGAP3 bind PI(3,4)P$_2$}~\cite{li2018}, and the direct exchange of 3.3b is the lumped form.
The feedback is not placed on the G protein, i.e.\ as a PIP3-accelerated dissociation of the heterotrimer (3.3g, available as an option), because \mk{M3-g-follows-stim}{G-protein activation provides a simple intracellular translation of the external gradient and reflects the local cAMP receptor occupancy}~\cite{xu2005} and not PIP3 (see Role), \mk{M3-gbg-pi3k-indep}{the cAMP-induced change of the G$\beta\gamma$ mobility and domain formation are independent of PI3K activity}~\cite{vanhemert2010} (mobility, not dissociation, was measured), and \mk{M3-sb-pip3-indep}{symmetry breaking of Ras requires G$\alpha_2$ and G$\beta\gamma$ but not the PIP3, cGMP, TorC2 and PLA2 pathways}~\cite{kortholt2013}.
A second positive feedback may act within the Ras network itself: \mk{M3-gef-pfb}{the co-localisation of RasGEFX and RasGEFB with Ras-GTP-enriched domains suggests that their recruitment depends on Ras-GTP, and hence a positive feedback between these GEFs and Ras-GTP}~\cite{iwamoto2025}. It is not modelled, so the excitability of the model runs through 3.3b only, whereas the cell is excitable also without PIP3 (see Role).
\item[3.3c$_\mathrm{R}$--e] \textbf{Brake: delayed negative feedback read from RasG-GTP; the mechanism is a candidate, the molecule is not identified.}
RasG-GTP activates the brake substrate B (3.3c$_\mathrm{R}$), B$^*$ decays with $\tau=10$\,s (3.3d) and accelerates the hydrolysis of RasG-GTP (3.3e).
\emph{Purpose.} The brake supplies the delayed negative feedback that excitability requires and terminates excursions that the stimulus did not cause (see Role and Adaptation and brake).
\emph{Reader.} The feedback is read from RasG-GTP and not from PIP3 or a PKB, because Ras remains excitable, and its response to cAMP is unchanged, without the PIP3, TorC2, PLA2 and sGC pathways (see Role); a feedback that terminates a Ras excursion therefore cannot require them.
Placing the feedback on Ras follows the published model, but not its form: \mk{M3-fuk-nfb}{that model places the positive and the delayed negative feedback on Ras, and its GAP recruitment to the membrane is promoted by Ras-GDP and opposed by Ras-GTP}~\cite{fukushima2019}; a GAP activity induced by RasG-GTP, as in 3.3c$_\mathrm{R}$, is the simpler form chosen here.
\emph{Candidate molecule.} A RasGAP that is activated through Ras exists: \mk{M3-c2gap-ras}{membrane translocation and activation of the RasGAP C2GAP1 require Ras on the membrane, i.e.\ a negative feedback with a buffer node}~\cite{xu2022} (the primary study is cited there and is not among the sources used here).
In the sources, however, it has a different role from the brake: \mk{M3-c2gap-adapt}{C2GAP1 controls both the basal Ras activity and the GPCR-mediated Ras adaptation}~\cite{xu2021c2gap}, whereas the brake is off at rest and does not set the adapted level, and whether C2GAP1 acts on RasG in particular is not shown.
C2GAP1 is therefore a candidate for the molecule, not evidence for the role; that a Ras-read GAP terminates spontaneous excursions is an assumption of the model.
\emph{Time constant.} $\tau$ is set by the refractory period: a recovery half-time of about 7\,s (see Adaptation and brake) corresponds to $\tau=t_{1/2}/\ln 2\approx10$\,s, assuming that the recovery is governed by the decay of B$^*$.
The half-time was measured with an F-actin reporter after mechanical stimulation, not for Ras, and since $\tau_I$ has the same value, the measurement does not separate the brake from the inhibitor.
\emph{Range (open contradiction).} The purpose and the candidate are local: a spontaneous excursion is a local event, and \mk{M3-c2gap-local}{C2GAP1 is a locally recruited RasGAP}~\cite{xu2021c2gap}. In the model, however, B and B$^*$ have the cytosolic diffusion constant of $20\,\mu\mathrm{m^2/s}$, so that within $\tau=10$\,s B$^*$ spreads over $\sqrt{D\tau}\approx14\,\mu$m, about one cell length, and the brake acts globally.
The global value was set for the earlier PIP3 reader, after a local brake ($D=0.4\,\mu\mathrm{m^2/s}$) had reduced the information about the gradient direction at the Ras stage to less than half of that without a brake; a local brake is strongest where the signal is strongest and thereby cancels the contrast, whereas a global brake subtracts the same amount everywhere (argued, not measured).
A local brake would terminate an excursion where it occurs but, by that comparison, costs gradient contrast; whether this also holds for the RasG reader has not been tested. The range is therefore a design choice that tuning has to decide (\texttt{DICTY\_M3\_DBRAKE}).
\emph{Alternative: PKB.} \mk{M3-miao-ext}{The negative regulation by the PKBs extends to other Ras proteins and to PI3K, which makes them candidates for the delayed negative feedback of the excitable network}~\cite{miao2017}, but \mk{M3-miao-mech}{they are suggested to act indirectly, by inhibiting the RasGEF Aimless and by activating PI5K}~\cite{miao2017}, and the documented PKB feedback spares RasG: \mk{M3-sca1-rasc}{the Sca1 complex regulates RasC activity while RasG activation is unaffected}~\cite{charest2010}. That edge is modelled in Module~8 (8.2p/8.2q/8.2s). The PKB reader (3.3c$'$) and the PIP3 reader (3.3c) remain available as options.
\emph{Status.} The brake is a proposed term. Its half-activation and strength are design targets that have not yet been tuned; tuning has to show that the rest state is the low branch, that the adapted plateau is not shifted by the brake, that no second excursion follows a step, and that the module resets as measured: \mk{M3-reset}{after removal of cAMP the network returns to its prestimulus state in less than 1\,min}~\cite{xu2022}.
\end{description}

\paragraph{Constants.}
The rate constants of the LEGI arms (3.1--3.5) are constrained by dynamical targets rather than by single measurements: an excitor time constant of 2\,s, an inhibitor time constant of 10\,s, both from the fit of the measured Ras response, and near-perfect adaptation of the plateau to a sustained step.
When $\tau_I$ was shortened from 30\,s to 10\,s, $k_{I,\mathrm{on}}$ and $k_{I,\mathrm{b}}$ were multiplied by the same factor $0.1/0.033=3.03$ as $k_{I,\mathrm{off}}$, so that all activated GAP fractions, and with them $k_\mathrm{cat}$, are unchanged and only the time constant moves; the fit tied the rates in the same way, as \mk{M3-tau-pair}{the activation rates were set to one tenth of the deactivation rates}~\cite[Suppl.\ Information]{takeda2012}. This change has not yet been recalibrated against the downstream modules.
The pool sizes were increased to reduce shot noise, and the rate constants were rescaled so that only the products $k_{E,\mathrm{on}}[\mathrm{GEF_R}]_\mathrm{tot}$, $k_{R,\mathrm{on}}[\mathrm{GEF_R}]_\mathrm{tot}$ and $k_{I,\mathrm{on}}[\mathrm{GAP}]_\mathrm{tot}$ enter the dynamics; the pool enlargement in Module~2 was paired with rescaling of $k_{E,\mathrm{on}}$ and $k_{I,\mathrm{on}}$ so that the activated fractions and time constants do not change.
Unpaired, the enlargement makes the stage worse, since the activated fraction of the excitor changes with the G$\beta\gamma$ level.
$K_\mathrm{NF1}$ is taken from the mammalian catalytic domain because no \emph{Dictyostelium} kinetics exist; the $k_\mathrm{cat}$ of the model is derived from the distance to the zero-order switch and is 11 times the mammalian value, so the active-GAP pool stands for more enzyme than is measured.
The feedback gain $k_\mathrm{PIP3,R}$ cannot be raised as a means to improve gradient contrast: because the feedback acts on the level, it lifts the resting baseline as much as the driven response, and a contrast is a ratio; above a critical gain the loop is bistable and stops resetting.
Of the brake constants only the decay has an anchor (the refractory period); the activation is derived from a half-activation at 6400 RasG-GTP molecules per voxel, the driven level recorded in the code, and the removal rate from a target strength (0.6 of the LEGI GAP arm at the driven state), so that the brake is off at rest and engages at the peak.
The brake and the GAP have a cytosolic diffusion constant of $20\,\mu\mathrm{m^2/s}$, which is appropriate for a cytosolic GAP but not for the membrane-associated brake B$^*$ of the code; it makes the brake global rather than local, which contradicts its local purpose and its local candidate molecule (see 3.3c$_\mathrm{R}$--e, Range) and is an open point of the present parameterisation.
The saturable GAP is discussed in Table~\ref{tab:hill}.
